%% kapa: exam=1 round=32 year=2021 period=2 form=A subject=accounting subject_ko=회계학 seq=30 number=70
%% source: Q-Net 감정평가사 기출문제 (cst00309) artlSeq=5208397 file=extracted/1008/5208397/A형/2021년도 제32회 감정평가사 1차 2교시 A형.pdf
%% conversion: kordoc 4.24.51 -> tools/gen1.py (auto); review=proofread-sonnet; flags=none
\begin{question}{70}
(주)감평은 20×1년 2월 초 영업을 개시하여 2년간 제품보증 조건으로 건조기(대당 판매가격 ￦100)를 판매하고 있다. 20×1년 1,500대, 20×2년 4,000대의 건조기를 판매하였으며, 동종업계의 과거 경험에 따라 판매수량 대비 평균 3\%의 보증요청이 있을 것으로 추정되고 보증비용은 대당 평균 ￦20이 소요된다. 당사가 제공하는 보증은 확신유형의 보증이며 연도별 보증이행 현황은 다음과 같다.

{\small\centering\begin{tabular}{|c|c|c|}\hline
구분 & 20×1년 & 20×2년 \\ \hline
20×1년 판매분 & 5대 & 15대 \\ \hline
20×2년 판매분 & & 30대 \\ \hline
\end{tabular}\par}
\vspace{4pt}
20×2년 말 보증손실충당부채는? (단, 보증요청의 발생가능성이 높고 금액은 신뢰성 있게 측정되었다. 충당부채의 현재가치요소는 고려하지 않는다.)
\begin{choices}
  \item ￦800
  \item ￦1,000
  \item ￦1,200
  \item ￦1,800
  \item ￦2,300
\end{choices}
\end{question}
\answer{5}
